\section{Further research questions and concrete next targets}
\label{sec:questions}

The following programme emphasizes exact identities and structural
reduction. Its first two items sharpen source questions partially
answered here; the subsequent items follow from the proved formulas.
Numerical experiments may guide these questions, but their desired
outputs are identities, rank theorems, and analytic proofs.

\subsection{The first \texorpdfstring{$q$}{q} moments for unequal even orders}

\begin{conjecture}[The optimal consecutive prefix]\label{conj:firstq}
For all even $2\le p<q$, the square matrix
\[
 (a_{jk})_{q/2\le j\le q-1,\ 0\le k<q/2}
\]
is nonsingular over $\Q$. Equivalently, the first $q$ centered
moments have the full formal rank $q$.
\end{conjecture}

The article proves a larger finite prefix for every pair and the
conjectured prefix on the diagonal. Exact calculations verify the
unequal assertion through $q=24$, but this is a finite test.
A useful next proof would replace the general exponential-polynomial
zero count by a sharper property of the restricted even-polynomial
space in \eqref{tr:operator}. An explicit determinant formula would
also explain the nontrivial rational numerators already visible for
$(2,4)$ and $(2,6)$. A sign pattern alone would be insufficient unless
it proves nonvanishing for every order.

\subsection{Can the mixed fifth coordinate be eliminated?}

Theorem~\ref{q5:main} singles out the holomorphic quantity $J'(1)$.
The first unresolved question is whether it has an exact reduction
to the height-one coefficients and ordinary zeta/Stieltjes constants
already present in that theorem. It is not appropriate to declare a
new independent constant merely because it appears in one reduction.
Conversely, a small numerical integer relation would not establish
its reducibility.

A more algebraic version asks for a scalar polynomial in $\psi$
whose coordinate finite part contains only a prescribed highest-depth
coefficient and ordinary constants. The scalar Gamma--Appell theorem
controls the nondecaying residue polynomial at infinity. Its explicit
defects, starting at degree four, show why this statement alone does
not settle all convergent lower-depth contributions. A promising
next target is to classify the residue defects modulo discrete
differences whose logarithmic sums reduce by finite harmonic algebra.

\subsection{An all-degree classification of the mixed coordinates}

The differential sequence $\CM_m$ has exactly one height-one residue
word. The raw scalar power $\psi^m$ introduces products involving
higher polygammas. At degree five the additional coordinate is represented
by $\FP\int\psi(\psi')^2$. At degree six and beyond, one can ask:
which differential monomials survive integration by parts after every
finite endpoint constant is retained?

A precise project is to grade monomials by
$\operatorname{wt}(\psi^{(r)})=r+1$, quotient them by total
derivatives, and compute the finite-dimensional residue-coordinate
image at each weight. This must keep endpoint evaluation as part of
the map: discarding a total derivative without its finite constant
would produce a different problem. The desired result is a recursion
for the surviving spectral words, with a proved coordinate count,
rather than a list of fitted numerical bases.

\subsection{Shifted and rational-cutoff Gamma normalizations}

Replacing the base point $1$ by a parameter $a$ naturally produces
$\Gamma(a-t)\Gamma(x)/(\Gamma(a)\Gamma(x-t))$ and differences
of Hurwitz polygammas. The special base point in this article makes
the origin residue exactly linear in $t$, which is responsible for
the ordinary integrals at every degree $m\ge2$. For general $a$, a
new subtraction is required.

Find the canonical joint subtraction in $x$ and $t$, prove the
resulting integral bridge, and express it in the ordered
multiple-Hurwitz Laurent coordinates already developed in the
incoming Barnes report. Rational $a$ should connect the resulting
formulas with congruence-restricted harmonic sums and cyclotomic
polylogarithms. The question is the exact compatibility of the
normalizations, not just formal substitution in an unshifted formula.

\subsection{Four tails and higher products}

For products $\prod_{i=1}^k(\zeta(p_i)-H_n^{(p_i)})$ with even
$p_i$ and even $k$, the centered even-power asymptotics persist.
Finite summation by
parts produces higher-depth convergent zeta coordinates. Determine a
finite formal coordinate system and prove an effective initial-moment
generation theorem for even $k\ge4$. Odd numbers of tails instead
produce odd centered powers and can have poles at the target spectral
values; their finite-part moments constitute a separate extension.
The two-tail proof suggests a route
through differential operators on products of $\coth(t/2)$, but
higher pole multiplicities and overlaps of coordinate columns must
be handled explicitly. Merely observing stabilization in low-order
matrices would not prove a uniform theorem.

\subsection{Exact primitives at algebraic or rational arguments}

The functions $\mathfrak P_m$ provide explicit normally convergent
logarithmic primitives on a slit plane. Evaluate their values at
rational arguments and their reflection or multiplication combinations.
One may seek reductions to Barnes multiple-Gamma functions, multiple
Hurwitz derivatives, or cyclotomic polylogarithms. The removable values
\eqref{eq:Cnegative} and the shift identity \eqref{eq:Cshift} supply
exact boundary conditions. Any proposed primitive identity should
match both its derivative and one normalization value; matching
derivatives alone leaves an undetermined constant.

\subsection{Higher Laurent jets and their functional relations}

Equation~\eqref{eq:alllaurent} supplies every regular Laurent
coefficient through a fixed logarithmic Mellin kernel. Determine
which identities among those logarithmic moments arise from
polylogarithmic functional equations, shuffle products, and the
spectral primitive \eqref{eq:spectral_primitive}. A particularly
concrete question is whether selected mixed derivatives in the
sixth-power calculation can be expressed using only finitely many
height-one jets plus a small collection of higher-depth germs.

\subsection{Proof certificates and unresolved ladder identities}

The finite harmonic identities in the fifth-power ledger are suitable
for independent formal verification: only exact quasi-shuffle
multiplication, the prefix rule, explicit low-weight reductions, and
polynomial constant-term extraction are needed. The analytic layer
can be separated into normal convergence, the contour remainder,
and the Mellin continuation lemmas. This provides a manageable route
to integrating proof certificates into ProveIt.

The manuscript's remaining $S_6$ and revised $S_8$ ladder candidates
continue to deserve exact functional-equation research. Their
existing high-precision certificates should guide the search for
an identity or an obstruction, while remaining explicitly distinct
from a proof of equality. No result of this article changes their
status.
